% Bài: L12C1 Bài 2. Nội năng. Định luật I của nhiệt động lực học · Lý thuyết về nội năng và các cách biến đổi nội năng
% ID: L12C1B2-01-DS
% Mức: NB
% ID: L12C1B2-01-DS


% Mức: NB
% ID: L12C1B2-76-TN
% ID: L12C1B2-76-TN
% Mức: NB
% ID: L12C1B3-118-TN
% ID: L12C1B3-118-TN
% Mức: NB
% ID: L12C1B2-75-TN
%=== Câu 84 ===%
\dangbt{Lý thuyết về nội năng và các cách biến đổi nội năng}
\begin{ex}
% ID: L12C1B2-75-TN
% Mức: NB
	Phát biểu nào sau đây đúng?
	\choice 
	{Khi một vật tỏa nhiệt ra môi trường thì nội năng của vật tăng lên}
	{Độ biến thiên nội năng của một vật là độ biến thiên nhiệt độ của vật đó}
	{Nội năng là phần năng lượng vật nhận được hay mất đi trong quá trình truyền nhiệt}
	{\True Nội năng của vật phụ thuộc vào nhiệt độ và thể tích của vật}
	\loigiai{
		\begin{itemize}
			\item Phát biểu  Nội năng của vật phụ thuộc vào nhiệt độ và thể tích của vật đúng vì nội năng của một hệ là tổng động năng và thế năng tương tác của các phân tử cấu tạo nên hệ. Động năng phân tử phụ thuộc vào nhiệt độ, còn thế năng tương tác phân tử phụ thuộc vào khoảng cách giữa các phân tử (tức là phụ thuộc vào thể tích của vật). Do đó, nội năng của vật phụ thuộc vào nhiệt độ và thể tích,,.
			\item Các phát biểu khác sai vì: Khi một vật toả nhiệt ra môi trường thì nội năng của vật giảm xuống. Phần năng lượng nhiệt mà vật nhận được hay mất đi trong quá trình truyền nhiệt được gọi là nhiệt lượng.
		\end{itemize}
	}
\end{ex}

% Mức: NB
% ID: L12C1B2-77-TN
% ID: L12C1B2-77-TN
% Mức: NB
% ID: L12C1B3-119-TN
% ID: L12C1B3-119-TN
% Mức: NB
% ID: L12C1B2-76-TN
%=== Câu 85 ===%
\begin{ex}
% ID: L12C1B2-76-TN
% Mức: NB
	Phát biểu nào dưới đây nói về nhiệt lượng là \textbf{không} đúng?
	\choice 
	{\True Một vật lúc nào cũng có nội năng, do đó lúc nào cũng có nhiệt lượng}
	{Nhiệt lượng là số đo độ biến thiên nội năng của vật trong quá trình truyền nhiệt}
	{Nhiệt lượng không phải là nội năng}
	{Đơn vị của nhiệt lượng cũng là đơn vị của nội năng}
	\loigiai{
		Nội năng là tổng động năng và thế năng của các phân tử cấu tạo nên vật, trong khi nhiệt lượng là phần nội năng trao đổi trong quá trình truyền nhiệt. Một vật luôn có nội năng nhưng không phải lúc nào cũng có nhiệt lượng, vì nhiệt lượng chỉ xuất hiện khi có sự truyền nhiệt giữa các vật.
	}
\end{ex}

% Mức: NB
% ID: L12C1B2-81-TN
% ID: L12C1B2-81-TN
% Mức: NB
% ID: L12C1B3-120-TN
% ID: L12C1B3-120-TN
% Mức: NB
% ID: L12C1B2-77-TN
%=== Câu 89 ===%
\begin{ex}
% ID: L12C1B2-77-TN
% Mức: NB
	Phần năng lượng nhiệt mà vật này truyền cho vật kia hoặc vật này nhận từ vật kia gọi là
	\choice 
	{nội năng}
	{thế năng}
	{nhiệt độ}
	{\True nhiệt lượng}
	\loigiai{
		Phần năng lượng nhiệt truyền từ vật có nhiệt độ cao hơn sang vật có nhiệt độ thấp hơn (hay phần năng lượng nhiệt trao đổi giữa hai vật) trong quá trình truyền nhiệt được gọi là nhiệt lượng,.
	}
\end{ex}
\dangbt{Lý thuyết về nội năng và các cách biến đổi nội năng}
\begin{ex}
Nội năng của một vật là:
\choice
{A. Tổng động năng và thế năng của vật.}
{B. Tổng động năng của các phân tử cấu tạo nên vật.}
{C. Tổng thế năng của các phân tử cấu tạo nên vật.}
{\True D. Tổng động năng và thế năng của các phân tử cấu tạo nên vật.}
\loigiai{Nội năng của một vật là tổng động năng và thế năng của các phân tử cấu tạo nên vật. Đây là định nghĩa cơ bản về nội năng.}
\end{ex}

\dangbt{Lý thuyết về nội năng và các cách biến đổi nội năng}
\begin{ex}
Khi nén một lượng khí trong xi lanh, nhiệt độ của khí tăng lên. Hiện tượng này chủ yếu là do:
\choice
{A. Khí nhận nhiệt từ môi trường bên ngoài.}
{B. Các phân tử khí va chạm mạnh hơn với thành xi lanh.}
{\True C. Môi trường bên ngoài đã thực hiện công lên khí.}
{D. Khí đã truyền nhiệt cho môi trường bên ngoài.}
\loigiai{Khi nén khí, piston thực hiện công lên khí. Công này làm tăng động năng trung bình của các phân tử khí, dẫn đến tăng nhiệt độ của khí. Đây là một cách biến đổi nội năng bằng cách thực hiện công.}
\end{ex}

\dangbt{Lý thuyết về nội năng và các cách biến đổi nội năng}
\begin{ex}
Một khối kim loại được nung nóng trên bếp. Cách biến đổi nội năng của khối kim loại trong trường hợp này là:
\choice
{A. Thực hiện công.}
{B. Thực hiện công và truyền nhiệt.}
{\True C. Truyền nhiệt.}
{D. Không có sự biến đổi nội năng.}
\loigiai{Khi khối kim loại được nung nóng trên bếp, nhiệt từ bếp được truyền trực tiếp vào khối kim loại, làm tăng nội năng của nó. Đây là hình thức biến đổi nội năng bằng cách truyền nhiệt.}
\end{ex}

\dangbt{Lý thuyết về nội năng và các cách biến đổi nội năng}
\begin{ex}
Đại lượng nào sau đây \textbf{không} phải là một cách để biến đổi nội năng của một vật?
\choice
{A. Thực hiện công lên vật.}
{B. Vật thực hiện công lên môi trường.}
{C. Truyền nhiệt cho vật.}
{\True D. Thay đổi khối lượng của vật.}
\loigiai{Nội năng của một vật có thể biến đổi bằng cách thực hiện công (công cơ học) hoặc truyền nhiệt. Thay đổi khối lượng của vật không phải là một cách biến đổi nội năng của vật đó theo định nghĩa vật lý nhiệt động lực học.}
\end{ex}

\dangbt{Lý thuyết về nội năng và các cách biến đổi nội năng}
\begin{ex}
Khi nghiên cứu về nội năng và sự biến thiên nội năng của một vật, xét tính đúng/sai của các phát biểu sau:
\choiceTF
{Nội năng của một vật là tổng động năng và thế năng tương tác của các phân tử cấu tạo nên vật}
{Nội năng của một lượng khí lí tưởng xác định chỉ phụ thuộc vào nhiệt độ của khối khí}
{Khi cọ xát một miếng kim loại trên mặt bàn, nội năng của miếng kim loại tăng lên do nó nhận nhiệt lượng từ mặt bàn}
{Trong quá trình truyền nhiệt không có sự chuyển hóa năng lượng từ dạng này sang dạng khác mà chỉ có sự truyền nội năng từ vật này sang vật khác}
\loigiai{
\begin{itemizelm}
    \item a) Đúng. Theo định nghĩa, nội năng của vật là tổng động năng chuyển động nhiệt và thế năng tương tác của các phân tử cấu tạo nên vật.
    \item b) Đúng. Vì các phân tử khí lí tưởng coi như không tương tác với nhau (bỏ qua thế năng tương tác phân tử), nên nội năng của khối khí lí tưởng chỉ phụ thuộc vào nhiệt độ.
    \item c) Sai. Nội năng của miếng kim loại tăng lên là do quá trình thực hiện công (cơ năng chuyển hóa thành nội năng), không phải do truyền nhiệt.
    \item d) Đúng. Bản chất của quá trình truyền nhiệt là sự truyền trực tiếp nội năng từ vật có nhiệt độ cao sang vật có nhiệt độ thấp hơn, không có sự chuyển hóa năng lượng giữa các dạng khác nhau.
\end{itemizelm}
}
\end{ex}

\dangbt{Lý thuyết về nội năng và các cách biến đổi nội năng}
\begin{ex}
Xét các phát biểu sau về các cách làm biến đổi nội năng của một hệ nhiệt động lực học:
\choiceTF
{Có hai cách làm biến đổi nội năng của một vật là thực hiện công và truyền nhiệt}
{Nhiệt độ của một vật tăng lên chứng tỏ các phân tử cấu tạo nên vật chuyển động nhiệt nhanh hơn, do đó nội năng của vật tăng}
{Nhiệt lượng là một dạng năng lượng mà vật luôn chứa đựng bên trong tương tự như nội năng}
{Khi nén một khối khí trong xilanh kín được bọc cách nhiệt hoàn toàn, nội năng của khối khí tăng lên do hệ nhận công}
\loigiai{
\begin{itemizelm}
    \item a) Đúng. Nội năng có thể biến đổi bằng hai hình thức: thực hiện công hoặc truyền nhiệt.
    \item b) Đúng. Nhiệt độ đặc trưng cho động năng chuyển động nhiệt trung bình của các phân tử, nhiệt độ tăng thì động năng phân tử tăng dẫn đến nội năng tăng.
    \item c) Sai. Nhiệt lượng không phải là dạng năng lượng vật chứa sẵn mà là số đo độ biến thiên nội năng trong quá trình truyền nhiệt.
    \item d) Đúng. Vì xilanh cách nhiệt nên $Q = 0$, quá trình nén khí là quá trình nhận công ($A > 0$), theo nguyên lí I nhiệt động lực học $\Delta U = A > 0$, do đó nội năng của khối khí tăng.
\end{itemizelm}
}
\end{ex}

\dangbt{Lý thuyết về nội năng và các cách biến đổi nội năng}
\begin{ex}
Một khối khí chứa trong xi lanh có pít-tông đóng kín nhận một công $150\text{ J}$ do ngoại lực nén vào. Đồng thời, khối khí truyền ra môi trường xung quanh một nhiệt lượng $40\text{ J}$. Độ biến thiên nội năng của khối khí bằng bao nhiêu jun ($\text{J}$)?
\shortans{110}
\loigiai{
Theo nguyên lí thứ nhất của nhiệt động lực học:
$\Delta U = A + Q$.
Vì khối khí nhận công từ ngoại lực nên $A = 150\text{ J}$.
Khối khí truyền nhiệt lượng ra môi trường nên $Q = -40\text{ J}$.
Do đó, độ biến thiên nội năng của khối khí là:
$\Delta U = 150 + (-40) = 110\text{ J}$.
}
\end{ex}

\dangbt{Lý thuyết về nội năng và các cách biến đổi nội năng}
\begin{ex}
Người ta truyền cho chất khí trong một bình kín một nhiệt lượng $240\text{ J}$. Do bình dãn nở, chất khí thực hiện một công $95\text{ J}$ đẩy pít-tông dịch chuyển. Độ biến thiên nội năng của lượng khí này bằng bao nhiêu jun ($\text{J}$)?
\shortans{145}
\loigiai{
Theo nguyên lí thứ nhất của nhiệt động lực học:
$\Delta U = A + Q$.
Chất khí nhận nhiệt lượng nên $Q = 240\text{ J}$.
Chất khí thực hiện công nên $A = -95\text{ J}$.
Độ biến thiên nội năng của khối khí là:
$\Delta U = 240 - 95 = 145\text{ J}$.
}
\end{ex}

\dangbt{Lý thuyết về nội năng và các cách biến đổi nội năng}
\begin{ex}
Trong một chu trình biến đổi, một khối khí nhận nhiệt lượng $350\text{ J}$ và nội năng của khối khí tăng thêm $120\text{ J}$. Công mà khối khí đã thực hiện lên môi trường ngoài bằng bao nhiêu jun ($\text{J}$)?
\shortans{230}
\loigiai{
Theo nguyên lí thứ nhất của nhiệt động lực học:
$\Delta U = A + Q$.
Khối khí nhận nhiệt lượng nên $Q = 350\text{ J}$.
Nội năng tăng thêm nên $\Delta U = 120\text{ J}$.
Công mà khối khí nhận được là:
$A = \Delta U - Q = 120 - 350 = -230\text{ J}$.
Công mà khối khí thực hiện lên môi trường ngoài là:
$A' = -A = 230\text{ J}$.
}
\end{ex}

\dangbt{Lý thuyết về nội năng và các cách biến đổi nội năng}
\begin{ex}
Nêu định nghĩa nội năng và các cách làm biến đổi nội năng của một vật. Cho ví dụ minh họa cho mỗi cách.
\loigiai{
* Định nghĩa nội năng: Nội năng của một vật là tổng động năng và thế năng của các phân tử cấu tạo nên vật. Nội năng phụ thuộc vào nhiệt độ và thể tích của vật.
* Các cách làm biến đổi nội năng:
    1. Thực hiện công: Khi có ngoại lực tác dụng lên vật làm vật dịch chuyển hoặc biến dạng, nội năng của vật có thể thay đổi.
        * Ví dụ: Dùng tay cọ xát mạnh hai lòng bàn tay vào nhau, lòng bàn tay nóng lên, nội năng của lòng bàn tay tăng.
    2. Truyền nhiệt: Khi có sự chênh lệch nhiệt độ giữa vật và môi trường xung quanh, nhiệt lượng có thể truyền từ nơi có nhiệt độ cao đến nơi có nhiệt độ thấp, làm thay đổi nội năng của vật.
        * Ví dụ: Đun nóng một ấm nước trên bếp, nước trong ấm nóng lên, nội năng của nước tăng.
}
\end{ex}

\dangbt{Lý thuyết về nội năng và các cách biến đổi nội năng}
\begin{ex}
Phân biệt sự khác nhau cơ bản giữa hai cách làm biến đổi nội năng: thực hiện công và truyền nhiệt. Nêu một ví dụ trong đó cả hai cách này cùng xảy ra đồng thời.
\loigiai{
* Sự khác nhau cơ bản:
    * Thực hiện công: Là quá trình biến đổi nội năng có liên quan đến sự trao đổi năng lượng dưới dạng cơ học (lực tác dụng làm vật dịch chuyển hoặc biến dạng). Trong quá trình này, có sự chuyển hóa từ cơ năng sang nội năng hoặc ngược lại.
    * Truyền nhiệt: Là quá trình biến đổi nội năng có liên quan đến sự trao đổi năng lượng dưới dạng nhiệt lượng do sự chênh lệch nhiệt độ. Trong quá trình này, không có sự chuyển hóa năng lượng từ dạng cơ học.
* Ví dụ cả hai cách cùng xảy ra đồng thời:
    * Khi bơm lốp xe đạp bằng bơm tay, không khí trong bơm bị nén (thực hiện công lên không khí), làm không khí nóng lên. Đồng thời, một phần nhiệt lượng từ không khí nóng truyền ra vỏ bơm và môi trường xung quanh (truyền nhiệt).
}
\end{ex}
